% GENERATED FILE. Edit canonical lessons, then run npm run generate:books.
% canonical-source-hash: fnv1a64:f44ec5c81b3f21d9
% canonical-lessons: GE-C259-aelter-als, GE-C259-seit, GE-C259-review

\chapter{Older Than, Since When — älter als, seit, vor}
\label{ch:ge-older-than-since-when}
\hlchaptermodality{259}

\begin{chapteropening}
\textbf{By the end of this chapter:} \emph{I can compare two things with -er als and so … wie, and say how long with seit and how long ago with vor.}

Complete GE-C259-review: compare two places or people, and say since when and how long ago, each with the right case.
\end{chapteropening}

\section[älter als, so alt wie]{älter als — older than}

\label{lesson:GE-C259-aelter-als}



\begin{quote}
Say \textquotedblleft{}old\textquotedblright{} and \textquotedblleft{}fast\textquotedblright{}. (\emph{Alt. Schnell.})

And remember \textbf{als}, from chapter two hundred and fifty-two: \textquotedblleft{}when\textquotedblright{}, for one time in the past. It has a second job.
\end{quote}

\begin{grammarlens}[title={Grammar Lens: -er, then als}]
To say \textquotedblleft{}more\textquotedblright{}, German adds \textbf{-er} to the adjective, as English does with \textquotedblleft{}faster\textquotedblright{}. \textquotedblleft{}Than\textquotedblright{} is \textbf{als}:

\begin{quote}
\emph{Der Zug ist schneller als das Auto.} — The train is faster than the car. \emph{Mein Bruder ist älter als ich.} — My brother is older than me.
\end{quote}

Many short adjectives with \emph{a}, \emph{o} or \emph{u} take an umlaut as well: \emph{alt}, \textbf{älter}; \emph{groß}, \textbf{größer}, bigger. Others do not: \emph{klein}, \emph{kleiner}; \emph{billig}, \emph{billiger}.

A few are irregular, as in English: \emph{gut}, \textbf{besser}; \emph{viel}, \textbf{mehr}; and \emph{gern}, \textbf{lieber}, which you already have.

After \emph{als}, a pronoun keeps its subject form: \emph{älter als ich}, \emph{als du}. In some regions and in casual speech you may hear \emph{älter wie ich}; in standard German it is \emph{als}.
\end{grammarlens}

\begin{culture}
For \textquotedblleft{}as … as\textquotedblright{}, German says \textbf{so … wie}, with the plain adjective in the middle:

\begin{quote}
\emph{Sie ist so alt wie ich.} — She is as old as me. \emph{Das Auto ist nicht so schnell wie der Zug.} — The car is not as fast as the train.
\end{quote}

\emph{Wie} is for the same, and \emph{als} is for more. \emph{So alt wie}, \emph{älter als}.
\end{culture}

\begin{cousinweb}
\textbf{alt} $\leftarrow$ Old High German \emph{alt} $\leftarrow$ Proto-Germanic \emph{*aldaz} (EtymDB 2.1), the parent of English \emph{old}. EtymDB does not list the comparative, so this part is the usual account: the old comparative ending had an \emph{i} in it, and that \emph{i} pulled the vowel forward, \emph{a} to \emph{ä}. English went through the same change and kept its trace in \emph{elder}, beside the newer \emph{older}.
\end{cousinweb}

\subsection*{Your turn}
Say these aloud:

\begin{itemize}
\raggedright
  \item \emph{alt, älter; klein, kleiner; gut, besser}
  \item \emph{Der Zug ist schneller als das Auto.}
  \item \emph{Mein Bruder ist älter als ich.}
  \item \emph{Sie ist so alt wie ich.}
  \item \emph{Ich nehme lieber Tee als Kaffee.} \emph{Lieber}, with \emph{als}
\end{itemize}

\begin{tcolorbox}[breakable,colback=teal!4,colframe=teal!35!black,title={Before you move on}]
\textquotedblleft{}Older than me\textquotedblright{}? (\emph{\textbf{Älter als ich.}}) \textquotedblleft{}As old as me\textquotedblright{}? (\emph{\textbf{So alt wie ich.}}) \textquotedblleft{}Better\textquotedblright{}? (\emph{\textbf{Besser.}}) Which English word kept the old umlaut of \emph{old}? (\emph{\textbf{Elder.}})
\end{tcolorbox}
\section[seit, vor]{seit and vor — since when, and how long ago}

\label{lesson:GE-C259-seit}



\begin{quote}
Say \textquotedblleft{}I am learning German\textquotedblright{}. (\emph{Ich lerne Deutsch.}) Now we add how long.
\end{quote}

\begin{grammarlens}[title={Grammar Lens: seit, with the present}]
English says \textquotedblleft{}I \textbf{have been} learning German for two years\textquotedblright{}. German keeps the present, because the learning is still going on:

\begin{quote}
\emph{Ich lerne} \textbf{seit} \emph{zwei Jahren Deutsch.}
\end{quote}

\textbf{seit} means \textquotedblleft{}since, for\textquotedblright{}, and it takes the dative: \emph{seit einem Jahr}, \emph{seit einer Stunde}, and in the plural an extra \emph{-n}: \emph{seit zwei Jahren}, \emph{seit drei Monaten}.

The question is \emph{seit wann?}, since when: \emph{Seit wann lernst du Deutsch?}

The trap for English speakers is the tense. \emph{Ich habe seit zwei Jahren Deutsch gelernt} forces the English shape onto German, and for something still going on it sounds wrong. If it is still true, keep the present.
\end{grammarlens}

\begin{culture}
For something over and done, German says \textbf{vor}, also with the dative, and it stands \textbf{before} the time, where English \textquotedblleft{}ago\textquotedblright{} comes after:

\begin{quote}
\emph{Vor zwei Jahren war ich in Berlin.} — Two years ago I was in Berlin.
\end{quote}

The time phrase takes the first place, so the verb comes next: \emph{vor zwei Jahren} \textbf{war} \emph{ich}.

And for a stretch of time that is finished, no little word is needed at all: \emph{Ich war drei Monate in Berlin}, I was in Berlin for three months.

So: still going on, \emph{seit} and the present. Over, \emph{vor} for when, or the bare time for how long.
\end{culture}

\begin{cousinweb}
\textbf{seit} $\leftarrow$ Middle High German \emph{sīt} (EtymDB 2.1). Further back it is usually traced to an Old High German word meaning \textquotedblleft{}later\textquotedblright{}, and linked with the first part of English \emph{since}; EtymDB does not record those steps. \textbf{vor} is usually traced to Old High German \emph{fora}, a cousin of English \emph{fore} in \emph{before}; EtymDB's entry for \emph{vor} does not settle it, so that too is given as the usual account.
\end{cousinweb}

\subsection*{Your turn}
Say these aloud:

\begin{itemize}
\raggedright
  \item \emph{Ich lerne seit zwei Jahren Deutsch.}
  \item \emph{seit einem Jahr, seit einer Stunde, seit drei Monaten}
  \item the question: \emph{Seit wann lernst du Deutsch?}
  \item \emph{Vor zwei Jahren war ich in Berlin.}
  \item \emph{Ich war drei Monate in Berlin.}
\end{itemize}

\begin{tcolorbox}[breakable,colback=teal!4,colframe=teal!35!black,title={Before you move on}]
\textquotedblleft{}I have been learning German for two years\textquotedblright{}? (\emph{\textbf{Ich lerne seit zwei Jahren Deutsch.}}) Which tense goes with \emph{seit} for something still true? (\textbf{The present.}) \textquotedblleft{}Two years ago\textquotedblright{}? (\emph{\textbf{Vor zwei Jahren.}})
\end{tcolorbox}
\section[(review)]{Review — then and now}

\label{lesson:GE-C259-review}



\begin{quote}
Say \textquotedblleft{}my brother is older than me\textquotedblright{}. (\emph{Mein Bruder ist älter als ich.})
\end{quote}

\begin{grammarlens}[title={What you've built}]
Two chapters, and an experience can be told from the inside: how it felt, what you hope, what it compares with, and when. \emph{Älter als}, \emph{so alt wie}. \emph{Seit} with the present for what is still true; \emph{vor} for what is over.

A short piece with \emph{vor}, \emph{als} and \emph{seit} in it:

\begin{quote}
\emph{Vor zwei Jahren war ich in Berlin. Der Zug war schneller als das Auto, und ich war froh. Seit einem Jahr lerne ich Deutsch. Ich hoffe, bald wieder nach Berlin zu fahren.}
\end{quote}

Listen for the verb after the time phrases: \emph{vor zwei Jahren} \textbf{war} \emph{ich}, \emph{seit einem Jahr} \textbf{lerne} \emph{ich}.
\end{grammarlens}

\subsection*{Your turn}
Say these aloud:

\begin{itemize}
\raggedright
  \item \emph{Vor zwei Jahren war ich in Berlin.}
  \item \emph{Seit einem Jahr lerne ich Deutsch.}
  \item \emph{Der Zug ist schneller als das Auto, aber nicht so schnell wie das Flugzeug.}
  \item \emph{Ich war froh, und ich hoffe, bald wieder nach Berlin zu fahren.}
  \item your own: how long you have been learning German, with \emph{seit}, and one thing from a year ago, with \emph{vor}
\end{itemize}

\begin{tcolorbox}[breakable,colback=teal!4,colframe=teal!35!black,title={Before you move on}]
\emph{Seit einem Jahr}: which tense comes with it, if it is still true? (\textbf{The present.}) \textquotedblleft{}A year ago\textquotedblright{}? (\emph{\textbf{Vor einem Jahr.}}) \textquotedblleft{}As fast as\textquotedblright{}? (\emph{\textbf{So schnell wie.}})
\end{tcolorbox}
